% 可编译组件示例。正式写作请替换内容；未用组件可注释保留。
\section{问题重述}
\subsection{问题背景}
此处填写问题背景。正文中文采用宋体，英文采用 Times New Roman；正文各段首行缩进两个字符。下面的三级标题紧接正文，不另加空白段。
\subsubsection{具体任务}
此处概括具体任务与约束条件。一级标题使用中文数字加顿号，二级与三级标题分别使用阿拉伯数字编号，例如 1.1 和 1.1.1。
\subsection{数据与任务说明}
说明数据来源、字段含义与预测或优化目标。示例仅用于展示排版，不对赛题做出未经验证的结论。

\section{问题分析与模型假设}
\subsection{问题分析}
此处填写分析与求解路线。文献引用的排序压缩示例为\citep{sample-c,sample-a,sample-b}；这里引用的是明确标注的格式占位条目，正式写作时必须替换为已核验来源。

本模板的排版、公式、图表占位、算法和附录代码示例由 OpenAI Codex 辅助生成与修改\citep{ai-codex}，不代表实际研究结果。正式写作时，应在实际采纳的 AI 辅助内容附近分别标注用途及人工核验情况。
\subsection{模型假设}
\begin{enumerate}
  \item 此处填写第一条有依据的假设。
  \item 此处填写第二条假设及其适用范围。
\end{enumerate}

\section{符号说明与表格示例}
表\ref{tab:symbols}是普通三线表；表\ref{tab:wide}演示长表题与自动换行列。全局列间距为零，因此用固定列宽或显式列间距组织表格。
\begin{table}[H]
  \centering
  \caption{主要符号说明}\label{tab:symbols}
  \begin{tabular}{>{\centering\arraybackslash}m{0.16\linewidth}
                  >{\raggedright\arraybackslash}m{0.54\linewidth}
                  >{\centering\arraybackslash}m{0.24\linewidth}}
    \toprule
    符号 & 含义 & 单位 \\
    \midrule
    $n$ & 样本数量 & 个 \\
    $x_i$ & 第 $i$ 个样本的观测值 & 按实际情况填写 \\
    $\hat y_i$ & 第 $i$ 个样本的预测值 & 按实际情况填写 \\
    \bottomrule
  \end{tabular}
\end{table}
\begin{table}[H]
  \centering
  \caption{用于检查多行表题居中、标签加粗和题注行距的示例：下面的数值仅作排版占位，不能用于论文中的实际结果报告}
  \label{tab:wide}
  \begin{tabularx}{0.94\linewidth}{>{\raggedright\arraybackslash}X
      >{\centering\arraybackslash}p{0.2\linewidth}
      >{\centering\arraybackslash}p{0.2\linewidth}}
    \toprule
    指标（示例数据） & 方案 A & 方案 B \\
    \midrule
    演示得分 & 0.82 & 0.91 \\
    对比说明：本列文字可根据可用宽度自动换行 & 待填写 & 待填写 \\
    \bottomrule
  \end{tabularx}
\end{table}

\section{数据预处理}
此处记录数据读取、缺失值处理、异常值处理和特征构造。输入与输出应使用项目相对路径，便于复现。

\section{模型建立与求解}
\subsection{编号公式}
所有独立的正式数学公式均采用有编号环境。线性回归的格式示例如式\eqref{eq:regression}：
\begin{equation}
  y_i=\beta_0+\beta_1x_i+\varepsilon_i.
  \label{eq:regression}
\end{equation}
多行推导可使用 \texttt{align}，每行保留右侧编号；式\eqref{eq:loss}与式\eqref{eq:update}演示对应的交叉引用。
\begin{align}
  L(\bm\beta)&=\frac{1}{n}\sum_{i=1}^{n}(y_i-\hat y_i)^2,\label{eq:loss}\\
  \bm\beta^{(t+1)}&=\bm\beta^{(t)}-\eta\nabla L(\bm\beta^{(t)}).\label{eq:update}
\end{align}
需要一个编号覆盖多行时，可在 \texttt{equation} 内使用 \texttt{aligned}，如式\eqref{eq:constraint}。
\begin{equation}\label{eq:constraint}
  \begin{aligned}
    \min_{\bm x}\quad & f(\bm x),\\
    \text{s.t.}\quad & g_j(\bm x)\le 0,\quad j=1,\ldots,m.
  \end{aligned}
\end{equation}
\subsection{核心流程伪代码}
算法\ref{alg:solver}演示带输入、输出、初始化、循环、判断、停止条件与返回值的流程，实际论文应替换为对应模型的求解逻辑。
\begin{algorithm}[H]
  \KwIn{数据 $D$，最大迭代次数 $T$，收敛阈值 $\epsilon$}
  \KwOut{当前最优解 $x^*$}
  初始化可行解 $x^*$，令 $t\leftarrow 0$，$\Delta\leftarrow+\infty$\;
  \While{$t<T$ 且 $\Delta>\epsilon$}{
    根据数据 $D$ 与当前解生成候选可行解 $x'$\;
    令 $\Delta\leftarrow|f(x')-f(x^*)|$\;
    \If{$f(x')<f(x^*)$}{更新 $x^*\leftarrow x'$\;}
    更新 $t\leftarrow t+1$\;
  }
  \Return{$x^*$}\;
  \caption{带收敛停止条件的迭代求解流程}\label{alg:solver}
\end{algorithm}
\subsection{单图与组合图}
图\ref{fig:single}是单图示例，按注释替换占位框即可插入实际图片。
\begin{figure}[H]
  \centering
  % \includegraphics[width=0.86\linewidth]{result.pdf}
  \fbox{\parbox[c][16mm][c]{0.78\linewidth}{\centering 此处放置实际结果图片}}
  \caption{单图题注示例}\label{fig:single}
\end{figure}
图\ref{fig:combined}保留子图和多行图题示例；主图题、子图题都居中，主图标签加粗。
\begin{figure}[H]
  \centering
  \begin{subfigure}[t]{0.46\linewidth}
    \centering
    \fbox{\parbox[c][12mm][c]{0.9\linewidth}{\centering 子图 A 占位}}
    \caption{方案 A}\label{fig:sub-a}
  \end{subfigure}\hfill
  \begin{subfigure}[t]{0.46\linewidth}
    \centering
    \fbox{\parbox[c][12mm][c]{0.9\linewidth}{\centering 子图 B 占位}}
    \caption{方案 B}\label{fig:sub-b}
  \end{subfigure}
  \caption{用于检查多行图题的组合图示例：替换为实际模型结果后，请核对图题是否居中、标签是否加粗以及图像是否被裁切}
  \label{fig:combined}
\end{figure}

\section{模型检验与评价}
此处填写误差评估、稳定性和敏感性分析，并说明模型局限。以上示例均不代表实际实验结论。

% 参考文献另起页；AI 声明在下一页独立排版，附录继续另起页。
\clearpage
\begin{thebibliography}{9}
  \bibitem{sample-a} 文献占位一：请替换为已经核验的作者、题名、出版信息和年份。此条仅用于展示数字引用。
  \bibitem{sample-b} 文献占位二：请替换为已经核验的真实来源。这里补充一段换行文字，以检查同一条目换行后的字号、基线距离与默认悬挂缩进。
  \bibitem{sample-c} 文献占位三：请替换为已经核验的真实来源，用于演示连续引用的排序与压缩。
  \bibitem{ai-codex} Codex，版本/型号：待核验，OpenAI，2026-10-23.（提交前按实际客户端和模型记录补全版本信息）
\end{thebibliography}
% 参考文献的下一页独立排版；正文在声明后另起页进入附录。
% 当前是据实填写的准备阶段记录；正式提交前须按最终论文更新。
\clearpage
\section*{AI工具使用声明}
\phantomsection
\addcontentsline{toc}{section}{AI工具使用声明}
截至 2026 年 10 月 23 日，本模板准备过程中使用了 OpenAI 提供的 Codex，辅助完成官方 Word 模板结构读取、LaTeX 模板与排版示例生成、样式调整、编译检查，以及 AI 使用规定的阅读与整理。具体模型及客户端版本待依据实际使用记录核验。模板中的公式、算法、图表占位和代码是排版示例，不能视作实际参赛成果。

正式论文所涉及的 AI 使用环节、采纳内容、人工修改与核验情况，应由参赛队在完成研究后逐项据实更新；本模板不预先声明尚未发生的独立建模或人工复核。详细记录另见支撑材料中的《AI工具使用详情说明》PDF。

% 后续提交时：保留并改写上述已使用 AI 的声明，补齐真实版本和使用日期；
% 对实际采纳的正文、图表、分析或程序分别标注 AI 贡献，并保留关键交互记录。
% 只有整项参赛作品确实从未使用任何 AI 工具时，才可改为规定的未使用声明。

\clearpage

% 不用 appendix 改为 A/B；附录沿用正文的中文一级标题与下级样式。
\section{附录}
附录的标题、普通文字、表格与正式公式继续沿用正文设置。代码框只额外设置等宽字体、行号、背景和边框。
\subsection*{问题一完整代码}
下面代码只演示相对路径、读取、预处理、参数、求解和关键输出；写作时须替换为实际可复现实验代码。
\begin{lstlisting}[language=Python]
# AI: Codex (OpenAI); version/model: verify before submission.
# Date: 2026-10-23; contribution: demonstration code.
# Human review: pending.
from pathlib import Path
import csv, json
def load_values(path):
    # Read the numeric input column using a relative path.
    with path.open(encoding="utf-8-sig", newline="") as stream:
        rows = csv.DictReader(stream)
        return [float(row["value"]) for row in rows
                if row.get("value", "").strip()]
def solve(values, offset=0.0):
    if not values:
        raise ValueError("The input has no valid values.")
    return sum(values) / len(values) + offset
if __name__ == "__main__":
    values = load_values(Path("data") / "input.csv")
    parameters = {"offset": 0.0}
    result = solve(values, **parameters)
    output = {"sample_count": len(values), "mean": result, "parameters": parameters, "note": "This intentionally long line demonstrates automatic wrapping inside the code frame."}
    print(json.dumps(output, ensure_ascii=False, indent=2))
\end{lstlisting}
\subsection*{问题二完整代码}
第二个代码框的行号重新从 1 开始；行号与代码字号相同。
\begin{lstlisting}[language=Python]
# AI: Codex (OpenAI); version/model: verify before submission.
# Date: 2026-10-23; demonstration code; human review pending.
def predict(values, intercept=0.0, slope=1.0):
    # Replace this demonstration with the actual model.
    return [intercept + slope * value for value in values]
print(predict([1.0, 2.0, 3.0]))
\end{lstlisting}
\subsection*{附录公式和表格}
附录中的正式公式仍应带编号，如式\eqref{eq:appendix-mean}。
\begin{equation}\label{eq:appendix-mean}
  \bar x=\frac{1}{n}\sum_{i=1}^{n}x_i.
\end{equation}
\begin{table}[H]
  \centering
  \caption{代码参数说明示例}\label{tab:appendix}
  \begin{tabular}{>{\raggedright\arraybackslash}p{0.25\linewidth}
                  >{\raggedright\arraybackslash}p{0.5\linewidth}}
    \toprule 参数 & 含义 \\
    \midrule \texttt{offset} & 均值的偏移参数，示例设为零 \\
    \texttt{slope} & 预测函数的斜率参数 \\
    \bottomrule
  \end{tabular}
\end{table}
